\section{Introduction}

Electronic medical records (EMRs) are commonly organized around healthcare providers.
Consequently, they primarily describe encounters that occur inside healthcare facilities and provide only intermittent snapshots of a patient's condition.
Patient-generated health data (PGHD) complements these snapshots with measurements, symptoms, and observations captured by patients through mobile applications, wearable devices, or manual entry.
Such data can improve clinicians' understanding of a patient's condition between visits and support more informed clinical encounters \cite{Cohen_Keller_Hayes_Dorr_Ash_Sittig_2016_PGHD_Integration_EHR,Omoloja_Vundavalli_2021_PGHD_Benefits_Challenges}.

Clinical integration of PGHD, however, is not merely an ingestion problem.
Prior reviews identify heterogeneous sources, workflow integration, unstable connectivity, data quality, privacy, and security as recurring obstacles \cite{Tiase2020-PGHD_EHR_Integration,Khatiwada_Yang_Lin_Blobel_2024_PGHD_Understanding}.
Sensor-derived PGHD is also high-frequency time-series data, making one remote transaction per measurement inefficient \cite{Haque_2024_Performance_Blockchain}.
A usable pipeline must therefore preserve measurements while a device is offline, control transmission volume, retain sufficient provenance, and prevent unauthorized parties from learning or modifying sensitive content.

DecMed is a decentralized EMR framework built on IOTA and the InterPlanetary File System (IPFS).
It provides patient-oriented identity and access-control mechanisms, but its original implementation remains provider-centric and has no path for continuous PGHD collection \cite{Purnama_2026}.
Directly placing raw PGHD on a distributed ledger would expose sensitive data and impose excessive transaction overhead.
Storing data only in a conventional backend, on the other hand, would weaken DecMed's decentralized audit and access-control model.

This work extends DecMed with a secure PGHD integration pipeline.
The main contributions are:

\begin{enumerate}
	\item an offline-first collection and batching mechanism for high-frequency PGHD from operating-system health APIs, device sensors, and manual input;
	\item a hybrid on-chain/off-chain architecture that combines client-side authenticated encryption, IPFS storage, IOTA metadata, and proxy re-encryption (PRE) for delegated access; and
	\item a provenance schema and an evaluation covering functional behavior, authorization, confidentiality, integrity, provenance completeness, workload capacity, and interruption recovery.
\end{enumerate}

The prototype is evaluated as an engineering proof of concept rather than a clinical system or production service-level benchmark.
Its purpose is to determine whether the selected mechanisms can preserve and protect PGHD across DecMed's complete patient-to-clinician data path.

\section{Background and Related Work}

\subsection{PGHD Integration Pipelines}

Zeng \textit{et al.} model PGHD integration as a pipeline containing acquisition, aggregation, and consumption functions \cite{Zeng2022-Standardized-PGD}.
Tiase \textit{et al.} found that existing integrations use portals, middleware, and data standards, but often remain pilot implementations whose clinical utility depends on workflow and visualization design \cite{Tiase2020-PGHD_EHR_Integration}.
These studies establish the functional stages required for PGHD integration, but they do not address DecMed's combination of decentralized identity, distributed metadata, and patient-controlled decryption.

\subsection{Decentralized Health Records and Provenance}

Akbulut and Tokdemir proposed an IOTA-based personal health record access-management system in which patients control access while the ledger supports integrity and traceability \cite{Akbulut_2020_iota_phr}.
DecMed develops a broader IOTA-based EMR framework with hierarchical access control \cite{Purnama_2026}.
Neither architecture originally provides an offline-first path for high-frequency PGHD.

Provenance is important because a value alone does not reveal who produced it, when it was measured, which device or application supplied it, or what processing occurred.
Healthcare provenance research consequently emphasizes traceable source and transformation information as a prerequisite for trustworthy reuse \cite{Ahmed_2023_Data_Provenance}.
This work records clinical provenance inside the encrypted PGHD payload while anchoring authorization, content references, and validity state in IOTA.

\subsection{Cryptographic Access Delegation}

PRE allows a semi-trusted proxy to transform encrypted key material for an authorized delegate without learning the underlying plaintext or the delegator's private key \cite{nunez2017proxy}.
This property is suitable for DecMed: a patient can grant a clinician time-bounded PGHD access without distributing the patient's private decryption key.
IPFS complements this mechanism with content-addressed off-chain storage, where a content identifier (CID) refers to encrypted content \cite{Benet_2014_ipfs}.
The proposed design combines these mechanisms instead of treating the ledger as bulk data storage.
